\documentclass[11pt,a4paper]{article}

\usepackage[utf8]{inputenc}
\usepackage[T1]{fontenc}
\usepackage[spanish]{babel}
\usepackage{lmodern}
\usepackage{geometry}
\usepackage{hyperref}
\usepackage{enumitem}
\usepackage{longtable}
\usepackage{array}
\usepackage{xcolor}
\usepackage{listings}

\geometry{margin=2.5cm}
\hypersetup{
    colorlinks=true,
    linkcolor=blue,
    urlcolor=blue,
    pdftitle={PCZT: desalineacion semantica entre recipient comprometido y recipient mostrado},
    pdfauthor={Codex}
}

\lstdefinestyle{plain}{
    basicstyle=\ttfamily\small,
    breaklines=true,
    frame=single,
    columns=fullflexible
}

\title{Analisis tecnico del bug de desalineacion semantica en PCZT}
\author{Documento interno de auditoria tecnica}
\date{\today}

\begin{document}

\maketitle

\tableofcontents
\newpage

\section{Resumen ejecutivo}

Este documento describe, de punta a punta, el proceso que seguimos para encontrar,
validar y convertir en PoC un bug semantico en el subsistema \texttt{PCZT}
(\textit{Partially Constructed Zcash Transaction}) dentro de \texttt{librustzcash}.

La conclusion central es la siguiente:

\begin{quote}
Es posible construir un PCZT valido cuyo output shielded comprometido paga al
recipient \texttt{B}, pero cuya metadata auxiliar hace que la capa de
reconstruccion de historial/salidas enviadas registre y exponga al recipient
\texttt{A}.
\end{quote}

En otras palabras:

\begin{itemize}[leftmargin=2em]
    \item \textbf{Recipient comprometido criptograficamente:} \texttt{B}
    \item \textbf{Recipient mostrado/almacenado en sent history:} \texttt{A}
    \item \textbf{La transaccion sigue siendo valida y se acepta:} si
\end{itemize}

Lo que \textbf{si} demostramos:

\begin{itemize}[leftmargin=2em]
    \item un flujo end-to-end dentro de \texttt{zcash\_client\_backend} donde:
    \begin{itemize}
        \item se crea un PCZT valido,
        \item se muta solo metadata de presentacion,
        \item se firman los spends,
        \item se extrae la transaccion final,
        \item y la capa de almacenamiento de salidas enviadas persiste un recipient diferente del comprometido.
    \end{itemize}
\end{itemize}

Lo que \textbf{no} demostramos todavia:

\begin{itemize}[leftmargin=2em]
    \item que una UI real de confirmacion de firma muestre \texttt{A} al usuario humano antes
    de firmar;
    \item que el mismo PoC este repetido para Orchard;
    \item que un wallet de produccion externo a este repo renderice esta metadata como
    fuente autoritativa en su pantalla de aprobacion.
\end{itemize}

\section{Contexto}

\subsection{Que es PCZT?}

\texttt{PCZT} es el formato de transaccion parcialmente construida de Zcash. La idea
es separar la construccion de una transaccion en roles:

\begin{itemize}[leftmargin=2em]
    \item \texttt{Creator}
    \item \texttt{Updater}
    \item \texttt{IoFinalizer}
    \item \texttt{Prover}
    \item \texttt{Signer}
    \item \texttt{Combiner}
    \item \texttt{Redactor}
    \item \texttt{SpendFinalizer}
    \item \texttt{TransactionExtractor}
    \item \texttt{Verifier}
\end{itemize}

Cada rol agrega, valida, poda o transforma informacion.

\subsection{Hipotesis inicial}

La hipotesis que queriamos probar era:

\begin{quote}
Un coordinador malicioso puede hacer que un participante vea semanticamente
``recipient A'' mientras la salida realmente comprometida por la transaccion
paga a ``recipient B''.
\end{quote}

La forma mas concreta de esto en el codigo era:

\begin{itemize}[leftmargin=2em]
    \item recipient comprometido en campos shielded del output;
    \item recipient de presentacion en campos auxiliares como:
    \begin{itemize}
        \item \texttt{user\_address}
        \item metadata propietaria \texttt{PROPRIETARY\_OUTPUT\_INFO}
        \item \texttt{PcztRecipient::External}
    \end{itemize}
\end{itemize}

\section{Preguntas que guiaron la investigacion}

Estas fueron las preguntas tecnicas que nos fuimos haciendo durante la busqueda:

\begin{enumerate}[leftmargin=2em]
    \item Cuales campos estan comprometidos criptograficamente y cuales son solo metadata?
    \item El recipient que ve un signer o un consumidor downstream sale del output real o de metadata?
    \item Hay algun chequeo generico que garantice:
    \[
    \texttt{user\_address} = \texttt{recipient comprometido}
    \]
    antes de firmar?
    \item El \texttt{Signer} de alto nivel valida la semantica de outputs o solo la validez minima para firmar?
    \item \texttt{extract\_and\_store\_transaction\_from\_pczt} reconstruye el recipient desde el note o desde metadata?
    \item Si mutamos solo \texttt{user\_address}, las firmas siguen siendo validas?
    \item Si las firmas siguen siendo validas, el historial guardado usa \texttt{A} o usa \texttt{B}?
\end{enumerate}

\section{Mapa de fuentes de verdad}

\subsection{Fuente de verdad criptografica}

Para outputs shielded, el recipient comprometido depende de:

\begin{itemize}[leftmargin=2em]
    \item en Sapling:
    \begin{itemize}
        \item \texttt{recipient}
        \item \texttt{value}
        \item \texttt{rseed}
        \item ciphertext
        \item commitment
    \end{itemize}
    \item en Orchard:
    \begin{itemize}
        \item \texttt{recipient}
        \item \texttt{value}
        \item \texttt{rho}
        \item \texttt{rseed}
        \item ciphertext
        \item commitment
    \end{itemize}
\end{itemize}

Esos campos alimentan:

\begin{itemize}[leftmargin=2em]
    \item el note comprometido,
    \item la recuperacion del output,
    \item la semantica on-chain real.
\end{itemize}

\subsection{Fuente de verdad de presentacion}

Para recipients externos, en el flujo PCZT aparecen dos piezas auxiliares:

\begin{itemize}[leftmargin=2em]
    \item \texttt{user\_address}
    \item \texttt{PROPRIETARY\_OUTPUT\_INFO} con \texttt{PcztRecipient::External}
\end{itemize}

Estas piezas \textbf{no} estan comprometidas en sighash ni en commitments.

\section{Trazado del flujo vulnerable}

\subsection{Creacion del PCZT}

En \texttt{zcash\_client\_backend/src/data\_api/wallet.rs}, el flujo
\texttt{create\_pczt\_from\_proposal} hace dos cosas importantes:

\begin{enumerate}[leftmargin=2em]
    \item toma el recipient semantico de la propuesta;
    \item lo traduce a metadata auxiliar:
    \begin{itemize}
        \item \texttt{user\_address}
        \item \texttt{PROPRIETARY\_OUTPUT\_INFO}
    \end{itemize}
\end{enumerate}

Para outputs externos, guarda:

\begin{lstlisting}[style=plain]
PcztRecipient::External
\end{lstlisting}

y tambien serializa la direccion de usuario.

\subsection{Firma}

El \texttt{Signer} de alto nivel:

\begin{itemize}[leftmargin=2em]
    \item calcula el \texttt{shielded\_sighash} desde \texttt{TransactionData<EffectsOnly>};
    \item firma los spends;
    \item no impone un chequeo output-side de consistencia entre:
    \begin{itemize}
        \item recipient comprometido,
        \item y \texttt{user\_address}.
    \end{itemize}
\end{itemize}

Esto es crucial: si mutamos \texttt{user\_address}, las firmas no cambian.

\subsection{Extraccion y almacenamiento}

En \texttt{extract\_and\_store\_transaction\_from\_pczt} ocurre el punto clave.

La funcion:

\begin{itemize}[leftmargin=2em]
    \item reconstruye el note desde los campos comprometidos;
    \item pero reconstruye el recipient externo para historial usando metadata:
    \begin{itemize}
        \item \texttt{PcztRecipient::External}
        \item \texttt{user\_address}
    \end{itemize}
\end{itemize}

En el helper \texttt{to\_sent\_transaction\_output(...)}:

\begin{lstlisting}[style=plain]
(PcztRecipient::External, Some(addr)) =>
    Recipient::External { recipient_address: addr, ... }
\end{lstlisting}

O sea: el recipient externo almacenado sale de metadata, no del recipient
comprometido del note.

\section{Hipotesis exacta del exploit}

La condicion necesaria y suficiente es:

\begin{enumerate}[leftmargin=2em]
    \item construir un PCZT valido que pague a \texttt{B};
    \item mutar solo \texttt{user\_address = A};
    \item dejar:
    \begin{itemize}
        \item recipient comprometido = \texttt{B}
        \item ciphertext / note / commitment = intactos
    \end{itemize}
    \item firmar normalmente;
    \item extraer/almacenar;
    \item observar si el recipient almacenado es \texttt{A}.
\end{enumerate}

\section{PoC ejecutable}

\subsection{Test agregado}

Se agrego el siguiente test:

\begin{lstlisting}[style=plain]
data_api::wallet::tests::
extract_and_store_transaction_from_pczt_can_store_metadata_recipient_distinct_from_committed_sapling_recipient
\end{lstlisting}

Archivo:

\begin{lstlisting}[style=plain]
zcash_client_backend/src/data_api/wallet.rs
\end{lstlisting}

\subsection{Que hace el test}

El test hace exactamente esto:

\begin{enumerate}[leftmargin=2em]
    \item crea una transaccion valida \texttt{transparent -> Sapling};
    \item recipient comprometido real:
    \begin{lstlisting}[style=plain]
    committed_recipient_addr = B
    \end{lstlisting}
    \item recipient falso de presentacion:
    \begin{lstlisting}[style=plain]
    fake_displayed_addr = A
    \end{lstlisting}
    \item construye el PCZT;
    \item muta solo:
    \begin{lstlisting}[style=plain]
    output.user_address = A
    \end{lstlisting}
    \item preserva:
    \begin{lstlisting}[style=plain]
    PcztRecipient::External
    \end{lstlisting}
    \item crea proofs;
    \item firma;
    \item finaliza spends transparentes;
    \item extrae la transaccion final;
    \item recupera el recipient comprometido via OVK del sender;
    \item persiste el resultado en un \texttt{MockWalletDb} instrumentado para capturar
    \texttt{SentTransactionOutput};
    \item compara el recipient recuperado con el recipient almacenado.
\end{enumerate}

\subsection{Comando de ejecucion}

\begin{lstlisting}[style=plain]
cargo test -p zcash_client_backend --features pczt \
  extract_and_store_transaction_from_pczt_can_store_metadata_recipient_distinct_from_committed_sapling_recipient \
  -- --nocapture
\end{lstlisting}

\subsection{Salida observada}

\begin{lstlisting}[style=plain]
running 1 test
test data_api::wallet::tests::
extract_and_store_transaction_from_pczt_can_store_metadata_recipient_distinct_from_committed_sapling_recipient
... ok
\end{lstlisting}

\subsection{Resultado semantico exacto}

El test prueba estas igualdades/desigualdades:

\begin{itemize}[leftmargin=2em]
    \item recipient comprometido recuperado del tx:
    \begin{lstlisting}[style=plain]
    recovered_recipient_addr == committed_recipient_addr == B
    \end{lstlisting}
    \item recipient mostrado/almacenado:
    \begin{lstlisting}[style=plain]
    displayed_recipient == fake_displayed_addr == A
    \end{lstlisting}
    \item mismatch:
    \begin{lstlisting}[style=plain]
    displayed_recipient != recovered_recipient_addr
    \end{lstlisting}
\end{itemize}

\subsection{Interpretacion}

Esto equivale a:

\begin{quote}
Committed recipient: B

Displayed/stored recipient: A

Transaction accepted: yes
\end{quote}

\section{Por que las firmas siguen siendo validas}

Porque el campo mutado, \texttt{user\_address}, es metadata auxiliar.

No entra en:

\begin{itemize}[leftmargin=2em]
    \item el \texttt{shielded\_sighash},
    \item el recipient comprometido del note,
    \item el ciphertext,
    \item el commitment.
\end{itemize}

Entonces:

\begin{itemize}[leftmargin=2em]
    \item el recipient real \texttt{B} queda criptograficamente intacto;
    \item la metadata \texttt{A} sigue viva para presentacion/almacenamiento;
    \item la verificacion criptografica no detecta el desalineamiento.
\end{itemize}

\section{Que parte del sistema consume la metadata falsificada}

La capa vulnerable no es la verificacion de la transaccion on-chain.

La capa vulnerable es la reconstruccion de \texttt{SentTransactionOutput}:

\begin{itemize}[leftmargin=2em]
    \item \texttt{Recipient::External} se construye desde metadata,
    \item el note/memo viene del output comprometido,
    \item el historial resultante mezcla ambas fuentes.
\end{itemize}

Eso convierte a la inconsistencia en una vulnerabilidad de integridad semantica
del historial/sumario.

\section{Preguntas que nos hicimos hasta llegar al PoC}

Estas fueron las preguntas mas utiles que guiaron el hallazgo:

\begin{enumerate}[leftmargin=2em]
    \item El recipient que ve el consumidor downstream sale del output real o de metadata?
    \item \texttt{user\_address} esta comprometido en sighash? No.
    \item Existe algun chequeo generico:
    \[
    \texttt{user\_address} = \texttt{recipient comprometido}
    \]
    antes de firmar? No.
    \item \texttt{extract\_and\_store\_transaction\_from\_pczt} usa metadata para
    clasificar recipients externos? Si.
    \item Entonces, si mutamos solo \texttt{user\_address}, se rompe la firma? No.
    \item Si no se rompe la firma, se persiste \texttt{A} o \texttt{B}? Se persiste
    \texttt{A}.
\end{enumerate}

\section{Que no demostramos todavia}

\subsection{Signer deception real frente a una UI humana}

No demostramos todavia que exista, dentro de este repo, una UI de confirmacion
de firma que muestre \texttt{user\_address = A} al usuario humano y le haga aprobar
mientras el output real paga a \texttt{B}.

Por eso la evaluacion final correcta es:

\begin{itemize}[leftmargin=2em]
    \item \textbf{Sent-history deception:} si, demostrado
    \item \textbf{Signer-confirmation deception:} incierto en este codebase
\end{itemize}

\subsection{Orchard}

No repetimos aun el mismo PoC completo para Orchard.

\section{Causa raiz}

La causa raiz no es una bug criptografica clasica.

Es una \textbf{fractura entre dos fuentes de verdad}:

\begin{itemize}[leftmargin=2em]
    \item fuente criptografica del recipient real;
    \item fuente auxiliar de presentacion del recipient externo.
\end{itemize}

Y esa fractura queda habilitada por tres propiedades simultaneas:

\begin{enumerate}[leftmargin=2em]
    \item la metadata de recipient externo es mutable;
    \item esa metadata no esta comprometida en la firma;
    \item la capa de reconstruccion/almacenado la usa como autoritativa.
\end{enumerate}

\section{Clasificacion tecnica}

La mejor forma de describir el bug es:

\begin{quote}
\textbf{semantic binding failure between committed shielded recipient and
stored/displayed external recipient metadata in PCZT-derived sent-history flows}
\end{quote}

\section{Severidad e interpretacion}

\subsection{Lo que si es}

\begin{itemize}[leftmargin=2em]
    \item bug de integridad semantica;
    \item coordinador malicioso puede sembrar metadata enganosa;
    \item el sistema puede almacenar y presentar esa metadata como recipient externo.
\end{itemize}

\subsection{Lo que no es}

\begin{itemize}[leftmargin=2em]
    \item no es forgery;
    \item no es bypass de ownership de spend;
    \item no es funds theft demostrada por si sola;
    \item no es inconsistencia de consenso.
\end{itemize}

\section{Recomendaciones de remediacion}

\begin{enumerate}[leftmargin=2em]
    \item No usar \texttt{user\_address} como fuente autoritativa de recipient externo si
    no fue verificada contra el recipient comprometido.
    \item Recomputar recipient de salida desde campos comprometidos siempre que sea posible.
    \item En el signer de alto nivel, agregar validacion explicita de coherencia output-side.
    \item Si se conserva metadata de presentacion, marcarla explicitamente como advisory y
    no mezclarla con objetos de ``recipient real''.
    \item Agregar el mismo PoC para Orchard.
\end{enumerate}

\section{Respuesta final}

\subsection{Is the mismatch fully demonstrated end-to-end?}

\textbf{PARTIALLY}

Explicacion exacta:

\begin{itemize}[leftmargin=2em]
    \item \textbf{Si}, esta completamente demostrado end-to-end para:
    \begin{itemize}
        \item construccion del PCZT,
        \item mutacion de metadata,
        \item firma,
        \item extraccion,
        \item almacenamiento de recipient mostrado/guardado.
    \end{itemize}
    \item \textbf{No}, no esta completamente demostrado para:
    \begin{itemize}
        \item una pantalla real de confirmacion de firma mostrada a un humano.
    \end{itemize}
\end{itemize}

La respuesta operacional mas precisa es:

\begin{quote}
El mismatch esta \textbf{plenamente demostrado} como bug de sent-history /
stored-recipient deception dentro del flujo PCZT de librustzcash.

Todavia esta \textbf{sin demostrar} como bug de signer-confirmation deception en una UI real.
\end{quote}

\end{document}
